\documentclass[10pt,a4paper]{amsart}
\usepackage[margin=19mm]{geometry}
\usepackage{amsmath,amssymb,mathtools}
\usepackage[scr=rsfs,cal=euler]{mathalfa}
\usepackage{xfrac}
\usepackage[unicode,hidelinks,bookmarks=false]{hyperref}
\newcommand{\ie}{i.e.\ }
\newcommand{\cf}{cf.\ }
\newcommand{\resp}{respectively\ }
\newcommand{\Subsection}{\subsection}
\providecommand{\nobreakdash}{\nobreak\hbox{-}}
\setlength{\emergencystretch}{2em}
\multlinegap0pt
\usepackage{bm}
\usepackage{bbm}
\usepackage{dsfont}
\usepackage{xspace}
\usepackage{cancel}
\usepackage{mathtools}
\usepackage{amsthm}
\makeatletter
\@ifundefined{thm}{\newtheorem{thm}{Thm}}{}
\@ifundefined{prop}{\newtheorem{prop}[thm]{Proposition}}{}
\makeatother
\newcommand{\sgn}{\mathrm{sgn}}
\title[Symbols for toric Eisenstein cocycles and arithmetic applica]{Symbols for toric Eisenstein cocycles and arithmetic applications}
\author{Peter Xu}
\date{Source: arXiv:2402.00294v5 (CC BY 4.0)}
\begin{document}
\maketitle

\noindent\emph{Source and licence.} Symbols for toric Eisenstein cocycles and arithmetic applications. Version arXiv:2402.00294v5 (CC BY 4.0), distributed under the Creative Commons Attribution 4.0 International licence, as recorded in the arXiv licence metadata for this exact version and shown on the abstract page https://arxiv.org/abs/2402.00294v5. Full rights notice: licenses/arxiv-2402.00294-CC-BY-4.0.txt.\par\smallskip

\noindent\textit{Editorial scope.} Counted unit: the Prop whose source label is \texttt{prop:gersrel} in the article (source0.tex lines 750--756, source label \texttt{prop:gersrel}), delivered together with the article's own complete original proof of it (source0.tex lines 757--864, one \texttt{proof} environment of 108 lines). No mathematics is written, repaired or re-derived: statement and proof are the article's own text line for line. Required context retained uncounted: eq:simpgen (lines 480--482). Every definition, display and label of those lines is retained unchanged. Omitted from this package and not redistributed: 1--479, 483--749, 865--2295. Nothing inside a retained excerpt is omitted or altered: every retained body is delivered exactly as the article's lines are written, so no omission receipt of any kind accompanies this package. Citations: the retained excerpts carry no citation command at all, so no bibliography entry of the article is transcribed and no reference is redistributed. Reference adaptation: none was needed and none was made; every reference inside the delivered unit resolves inside it, so no unresolved reference or citation is produced. Printed slips kept literal: no printed word, symbol, hypothesis or number of the counted statement or of its proof was altered. Layout and class: the article's own class is \texttt{amsart} with packages including graphicx, epsfig, mathabx, bbm, appendix, esvect, tikz-cd, fancyhdr, amsfonts, amsmath, amssymb, amsthm, parskip, times; the wrapper is the lane's pinned \texttt{amsart} editorial wrapper at 10pt on A4, which restates the theorem-like environments the retained text needs and adds the article's own macro definitions that the retained text needs (\texttt{\textbackslash sgn}), with extra wrapper packages (bm, bbm, dsfont, xspace, cancel, mathtools). The delivered numbering is the unit's own. Assets: no retained span contains a figure environment, a graphics-inclusion command, a PDF-page inclusion, a table or a TikZ picture; no image file is copied into this package, the package declares no asset dependency and no image file is redistributed with it. Licence: version arXiv:2402.00294v5 of this article is distributed under the Creative Commons Attribution 4.0 International licence, as recorded in the arXiv licence metadata for this exact version and shown on the abstract page https://arxiv.org/abs/2402.00294v5. A full rights notice is delivered at licenses/arxiv-2402.00294-CC-BY-4.0.txt. Characters: every retained excerpt is pure ASCII, so the delivered engine, XeLaTeX with its default Latin Modern text font, reports no missing character.
        \begin{equation}\label{eq:simpgen}
        [\Delta(m_1,\ldots, m_k)]=\sum_{i=1}^r [\Delta(m_1,\ldots, \hat{m_i},m,m_{i+1},\ldots, m_k)].
        \end{equation}

    \begin{prop} \label{prop:gersrel}
    For any integer $2\le r \le n$, each acyclic independent tuple $\underline{m}=(m_1,\ldots, m_n)$ of rays and ray $m$ lying on the great circle corresponding to the face spanned by $(m_1,\ldots, m_n)$ and sharing some hemisphere with all of them, the relation
        \begin{equation}
        \rho[\Delta(m_1,\ldots, m_k)]=\sum_{i=1}^r \rho[\Delta(m_1,\ldots, \hat{m_i},m,m_{i+1},\ldots, m_n)]
        \end{equation}
    coming from \eqref{eq:simpgen} holds. 
    \end{prop}

    \begin{proof}
        We may reduce to the case $r=n$ as follows: the claimed relation can be written as 
        \[
        \begin{pmatrix}m_1&\ldots & m_n\end{pmatrix}_* \{1-z_1,\ldots, 1-z_n\} = \sum_{i=1}^r \begin{pmatrix}m_1&\ldots& \hat{m_i} & m & \ldots   & m_k\end{pmatrix}_* \{1-z_1,\ldots, 1-z_n\}.
        \]
        The left-hand side factors as the cup product
        \[
        \begin{pmatrix}m_1&\ldots & m_r\end{pmatrix}_* \{1-z_1,\ldots, 1-z_r\} \smile \begin{pmatrix}m_{r+1}&\ldots & m_n\end{pmatrix}_* \{1-z_{r+1},\ldots, 1-z_n\} 
        \]
        and the right-hand side as 
        \[
        \left(\sum_{i=1}^r \begin{pmatrix}m_1&\ldots& \hat{m_i} & m & \ldots   & m_r\end{pmatrix}_* \{1-z_1,\ldots, 1-z_r\}\right)
        \smile \begin{pmatrix}m_{r+1}&\ldots & m_n\end{pmatrix}_* \{1-z_{r+1},\ldots, 1-z_n\} 
        \]
        so it suffices to prove that 
        \[
        \begin{pmatrix}m_1&\ldots & m_r\end{pmatrix}_* \{1-z_1,\ldots, 1-z_r\} = \sum_{i=1}^r \begin{pmatrix}m_1&\ldots& \hat{m_i} & m & \ldots   & m_r\end{pmatrix}_* \{1-z_1,\ldots, 1-z_r\}
        \]
        But this is the pullback of a top-rank stellar relation from the quotient $\mathbb{G}_m^n\twoheadrightarrow \mathbb{G}_m^n/G$, for $G$ the image of 
        \[
        \begin{pmatrix}m_{r+1}&\ldots & m_n\end{pmatrix}: \mathbb{G}_m^r \to \mathbb{G}_m^n.
        \]
        We therefore henceforth assume that $r=n$, and need to prove the relation 
        \[
        \begin{pmatrix}m_1&\ldots & m_n\end{pmatrix}_* \{1-z_1,\ldots, 1-z_n\} - \sum_{i=1}^n \begin{pmatrix}m_1&\ldots& \hat{m_i} & m & \ldots   & m_n\end{pmatrix}_* \{1-z_1,\ldots, 1-z_n\} =0.
        \]
        Taking the Bloch cycle description
        \[
        K_n^M(k(\mathbb{G}_m^n)) \hookrightarrow z^n(k(\mathbb{G}_m^n)\times \Box^n)/\partial z^{n+1}(k(\mathbb{G}_m^n)\times \Box^{n+1})
        \]
        we see that it suffices to show that 
        \begin{equation} \label{eq:finalstell}
        \sum_{i=1}^{n+1} (-1)^i \begin{pmatrix}m_1&\ldots& \hat{m_i} & \ldots   & m_{n+1}\end{pmatrix}_* \Gamma(1-z_1,\ldots, 1-z_n) \in \partial z^{n+1}(k(\mathbb{G}_m^n)\times \Box^{n+1}).
        \end{equation}
        for all acyclic tuples of rays $(m_1,\ldots, m_{n+1})$ such that each sub-$n$-tuple is full rank. Formally, \eqref{eq:finalstell} is the image under the cubical face maps of
        \begin{equation} \label{eq:finalcycle}
        \begin{pmatrix}m_1&\ldots & m_{n+1}\end{pmatrix}_* \Gamma(1-z_1,\ldots, 1-z_{n+1})\in \partial z^{n+1}(k(\mathbb{G}_m^n)\times \Box^{n+1})
        \end{equation}
        where the matrix denotes the map
        \begin{equation}\label{eq:n+1}
        \begin{pmatrix}m_1&\ldots & m_{n+1}\end{pmatrix}: \mathbb{G}_m^{n+1}\to \mathbb{G}_m^n.
        \end{equation}
        However, this seems not to use the acyclicity condition; what gives? The problem is that \eqref{eq:n+1} is not a finite map, so the pushed forward cycle may have improper intersection with the faces. Indeed, we claim that \eqref{eq:finalcycle} intersects all faces properly exactly when $(m_1,\ldots, m_{n+1})$ is acyclic. We will prove this on the level of cycles in 
        \[
        \mathbb{G}_m^n \times \Box^{n+1}
        \]
        since this certainly implies it for the restriction to the generic point. Write $M$ for the matrix $(m_1,\ldots, m_{n+1})$, and define the $(n+1)$-variable monomials
        \[
        p_j(z_1,\ldots, z_{n+1}) = \prod_{i=1}^{n+1} z_i^{M_{j,i}}
        \]
        whose exponents are the $j$th row of $M$. Then in $\mathbb{G}_m^n \times \Box^{n+1}$, the cycle \eqref{eq:finalcycle} can be described as the closure of the locus
        \[
        (p_1, p_2, \ldots, p_n, 1-z_1,\ldots, 1-z_{n+1}). 
        \]
        Its intersection with a codimension-$d$ cubical face is then indexed by a labelled subset $I\in \{1,\ldots, n+1\}$ of cardinality $d$, with each element $i$ of the subset labelled by $\ell_I(i)\in \{0,\infty\}$; this face is given by the intersection of the cycle with the locus
        \[
        \bigcap_{i\in I} \{ 1-z_i = \ell_I(i)\} 
        \]
        i.e. fixing each $z_i$ corresponding to the indexing set to be either $1$ or $\infty$. We wish to check, for each $I$, whether or not this intersection has the correct codimension, i.e. codimension $n+|I|$; the cycle \eqref{eq:finalcycle} meets all faces properly if and only if all codimensions are correct. 

        First consider the case where $I$ has at least one label of $0$; without loss of generality, we assume that it corresponds to $1\in I$, i.e. the relation $1-z_1=1\Leftrightarrow z_1=0$. The intersection of \eqref{eq:finalcycle} with $\{z_1=0\}$ is then the closure of the locus
        \begin{equation} \label{eq:minus1loc}
        (p_1(z_1=0),\ldots, p_n(z_1=0), 0, 1-z_2,\ldots, 1-z_n).
        \end{equation}
        where the notation indicates that we plug in $0$ in the $i$th place. Write $M'$ for the submatrix of $M$ given by deleting the first column; by assumption, it has full rank $n$. Viewing $M'$ as associated to a map
        \[
        \mathbb{G}_m^n \times \Box^n \to \mathbb{G}_m^n \times \Box^{n+1}
        \]
        by its natural action on the toric part, and by inclusion $\Box^n\hookrightarrow \Box^{n+1}$ in the last $n$ coordinates, we see that the locus \eqref{eq:minus1loc} is then the finite pushforward
        \begin{equation}\label{eq:minus1cyc}
        M'_* \Gamma(1-z_1,\ldots, 1-z_n).
        \end{equation}
        The intersection of \eqref{eq:finalcycle} with a face corresponding to $I$ is then the intersection of \eqref{eq:minus1cyc} with the face corresponding to the labelled set $I\setminus \{1\}$. But \eqref{eq:minus1cyc} is a finite pushforward of a cycle meeting all faces properly, so we conclude that \eqref{eq:finalcycle} meets the face corresponding to $I$ properly as well. Thus, any face with at least one label of $0$ always intersects \eqref{eq:finalcycle} properly; it therefore suffices to check the intersection with faces labelled only with $\infty$s. 

        We claim that if $(m_1,\ldots, m_{n+1})$ is acyclic, this intersection is always empty, and thus trivially proper.\footnote{When $(m_1,\ldots, m_{n+1})$ fails to be acyclic, the corresponding stellar simplicial relation should not hold, and thus \eqref{eq:finalcycle} must not meet all faces properly. The simplest example of what happens in this case: if $n=1$, $m_1=1$, and $m_2=-1$, then the locus \eqref{eq:finalcycle} is the closure of $(z_1z_2^{-1},1-z_1,1-z_2)$. Its intersection with the unique codimension-two face labelled with two $\infty$s should therefore be codimension $3$, i.e. zero-dimensional. However, the points of the curve parameterized by $(g,\infty,\infty)$ are in the closure for all $g\in \mathbb{G}_m$, since the point corresponding to each fixed $g$ is in the closure (as $t\to \infty$) of the curve $z_1=gt$, $z_2=t$ with free parameter $t$.} Note that the property of having empty intersection with the $\infty$-labelled faces is invariant under left multiplication of $M$ by elements of $\mathrm{GL}_n(\mathbb{Q})\cap M_n(\mathbb{Z})$. Via left multiplication by such a matrix, we can always turn the first $n$ columns of the matrix into scalar multiples of the standard basis. Thus, it suffices to check matrices of the form
        \[
        \begin{pmatrix}
        x_1 &  & & & y_1 \\
            & x_2    &  & & y_2 \\
            &      & \ldots &  & \ldots \\
            &        & & x_n   & y_n
        \end{pmatrix}
        \]
        where each of the $x_i$ and $y_i$ must be nonzero by the assumption that every $n$-by-$n$ submatrix of $M$ is full rank, and $\sgn(x_i)=\sgn(y_i)$ for at least one $i$ by the acyclicity assumption.

        In this case, suppose without loss of generality that $I=\{1,2,\ldots, k\}$ for some $k\le n$, with $\ell(j)=\infty$ for each $j\in I$. Assume now for the sake of contradiction that \eqref{eq:finalcycle} intersects the face corresponding to $I$ nontrivially; in particular, suppose it contains a point with $z_1=z_2=\ldots=z_k=\infty$ but 
        \[
        p_i=z_i^{x_i}z_{n+1}^{y_i} \in \mathbb{G}_m(\overline{k})
        \]
        for $\sgn(x_i)=\sgn(y_i)$. If $i\in \{1,2,\ldots, k\}$, this is impossible, since at this point we must have $z_{n+1}=0\Rightarrow 1-z_{n+1}=1$, which means our point fails to lie in the algebraic cube $\Box^{n+1}$. Otherwise, $x_j$ and $y_j$ have opposite signs for $j=1,2,\ldots, k$, meaning that $z_{n+1}=\infty$ as well at this point. But then
        \[
        p_i=z_i^{x_i}z_{n+1}^{y_i} \in \mathbb{G}_m(\overline{k})
        \]
        implies that $z_i=0\Rightarrow 1-z_{i}=1$, again a contradiction. We conclude the intersection with the face corresponding to $I$ is in fact empty, as desired.

        It remains only to show that the fundamental class of $S^{n-1}$ is sent to the generator
        \[
        \{-z_1,\ldots, -z_n\}.
        \]
        Indeed, we can decompose the fundamental class as a sum of simplicial orthants 
        \[
        [S^{n-1}]=\sum_{I\in \{\pm 1\}^n} [\Delta(I)]
        \]
        where $\Delta(I)$ is the simplex corresponding to $\sigma(I) (I_1\cdot e_1,\ldots, I_n\cdot e_n)$ where $\sigma(I)$ is an arbitrary even permutation if $(I_1\cdot e_1,\ldots, I_n\cdot e_n)$ is positively oriented, is an arbitrary odd permutation otherwise. Under $f_n$, this is sent to the sum
        \[
        \sum_{I\in \{\pm 1\}^n} \sigma \{ 1-z_1^{I_1},\ldots, 1-z_n^{I_n}\} = \left\{ \frac{1-z_1}{1-z_1^{-1}} ,\ldots \frac{1-z_n}{1-z_n^{-1}} \right\} = \{ -z_1,\ldots, -z_n\}.
        \]
    \end{proof}

\end{document}
